1つ面白い話をしよう.
まず,次の3つの漸化式を見てほしい.
\[
  a_{n+1}=a_n^2,
  \quad
  a_{n+1}=a_n^2-1,
  \quad
  a_{n+1}=a_n^2-2
\]
1つ目はきちんとここまで本書を読んできた諸君らならば,容易に解くことができるだろう.
実際,繰り返し代入すれば
\[
  a_n=a_1^{2^{n-1}}
\]
と分かる.\,$a_1>0$なら,両辺の対数を取って等比数列に帰着させることもできる.
一方,\,$a_1=0$や$a_1<0$では実数の対数を最初から取ることはできないが,元の漸化式そのものは定義されている.

2つ目はどうだろうか.\,$a_n$に指数がついているから,対数を使って処理したいが,\,$-1$の項があるためそのままではうまくいかない.
ただし,不動点が存在しないわけではない.
不動点$\alpha$は
\[
  \alpha=\alpha^2-1
  \quad\Longleftrightarrow\quad
  \alpha=\frac{1\pm\sqrt5}{2}
\]
である.
したがって,初項をどちらかの不動点に選べば,\,$a_n=\alpha$という定数解が得られる.
これは,後で扱う\sectionref{節：初項依存型}にもつながる考え方である.

しかし,不動点が見つかることと,一般の初項に対して漸化式がこれまでの型へ帰着できることは別である.
実際,\,$b_n=a_n-\alpha$と置くと,不動点の関係$\alpha^2-\alpha-1=0$より
\[
  b_{n+1}=b_n^2+2\alpha b_n
\]
となり,二次の項が残る.
つまり,不動点を引くだけでは等比型などの線形な漸化式にはならない.
少々意地悪であったかもしれないが,これは(\sectionref{節：漸化式の分類}で述べた意味において)本書でここまで扱った手法だけから一般の初項に対する閉じた一般項を得ることが難しい漸化式である.

3つ目も一見先の例のように一般項が得難い見た目をしている.
しかしながらこれはある特殊な条件の噛み合いによって偶然にも初等的な解法が存在する.

ここまで,等差型,等比型,階差型,階比型など,漸化式を一定の型に
帰着して解く方法を学んできた.しかし,右辺が少し複雑になると,
同じ道具が急に働かなくなる.

ここで扱うのは,一般の漸化式がいつでも解けるという話ではない.
一見すると解けそうにない漸化式でも,既知の数学的対象と一致する構造を持つことがある.
その一致を見抜くことで,初等的な型に還元できない漸化式を解いていく.
本章はそういう話である.